\documentclass[a4paper,12pt,twoside]{ctexart}
\usepackage[dvipsnames]{xcolor}
\usepackage{geometry}
\usepackage{tcolorbox}
\usepackage{tikz}
\usepackage{amssymb}
\usepackage{graphicx}

% 定义公函标准炭黑（K=92%，RGB约0.08）：权威稳重、高质量印刷感
\definecolor{inkblack}{rgb}{0.08,0.08,0.08}

\geometry{a4paper, left=3.18cm, right=3.18cm, top=2.54cm, bottom=2.54cm}

% 配置页脚，实现奇数页右下角，偶数页左下角的页码
\usepackage{fancyhdr}
\pagestyle{fancy}
\fancyhf{} % 清除所有页眉页脚的默认内容
\renewcommand{\headrulewidth}{0pt} % 移除页眉的横线

% 自定义页脚位置微调，完美还原绝对物理坐标（奇数页 x≈533.3，偶数页 x≈56.6，y≈788.5）
% 右下角（RO）：由于右侧文本域边缘位于 505pt，我们需要让它继续往右突破边界约 33.45pt
\fancyfoot[RO]{\makebox[0pt][l]{\hspace*{28.2pt}\raisebox{2.81pt}{\rmfamily\fontsize{10.5bp}{10.5bp}\selectfont\thepage}}}
% 左下角（LE）：左侧文本域边缘位于 90.14pt，我们需要让它向左突破边界移到 56.69pt（向左33.45pt）
\fancyfoot[LE]{\makebox[0pt][r]{\hspace*{33.45pt}\raisebox{2.81pt}{\rmfamily\fontsize{10.5bp}{10.5bp}\selectfont\thepage}}}

% 声明使用本地的方正大标宋字体
\newCJKfontfamily\fzdbs[Path=./]{FZDBS.ttf}
% 声明使用本地的方正小标宋字体（正文）
\newCJKfontfamily\fzssk[Path=./]{FZSSK.ttf}

% 全局文档颜色设定：公函炭黑
\AtBeginDocument{\color{inkblack}}

% 定义复合标题样式，严格匹配 PDF 中的字号与间距
% #1: 主标题（如：刑事（附带民事）自诉状） - 字号 22bp
% #2: 副标题（如：（侮辱案）） - 字号 18bp
\newcommand{\lawtitle}[2]{
    % 原始 PDF 主标题顶部 y 坐标为 110.7
    \vspace*{0.507cm} 
    \begin{center}
        {\fzdbs\fontsize{22bp}{22bp}\selectfont #1}\par
        \vspace{10.37bp} % 调整主标题与副标题间距，以匹配副标题顶部 y 坐标 (~144.4)
        {\fzdbs\fontsize{18bp}{18bp}\selectfont #2}\par
    \end{center}
    \vspace{18.4bp} % 标题下方的留白向下文中过渡
}

% 定义表框，严格对齐 PDF 中的表格宽度（467.72bp，对应原左 70.86，右 538.58）
\newtcolorbox{notesection}[1][]{
    width=467.72bp,      % 直接写出像素级宽度
    sharp corners,       % 直角边框
    colback=white,          % 背景白
    colframe=inkblack,     % 边框炭黑（公函品质）
    boxrule=0.5pt,         % 极细线条 0.5pt
    fontupper=\color{inkblack}\fontsize{10.5bp}{10.5bp}\selectfont, % 基础字体设置，不强制黑体
    enlarge left by=-19.28bp, % 修正左侧X坐标偏移对齐到70.86
    before=\vspace*{-13.39bp}\noindent, % 修正Y坐标偏移对齐到180.48
    after=\par,                          % 仅输出\par，消除默认后间距
    boxsep=0pt,          % 清除一切额外内白
    left=4bp,            % 严格内边距：左侧文字距黑线 4bp 达到绝对坐标 x=75.36
    top=4.69bp,          % 严格内边距：顶部文字距黑线 4.69bp 达到绝对坐标 y=184.3
    bottom=4.69bp,          % 严格内边距：底部文字距黑线 4.69bp，与顶部对称
    right=0pt,           % 严格右侧内边距
    #1
}

% 法律段落换行命令：等价于 \par + 修正 ctexart 额外行距的 -4.86bp 偏移
% 用于在严格 16bp 等高网格中保持每行精确对齐原档坐标
\newcommand{\lawpar}{\par\vspace*{-4.86bp}}

% 定义区块副标题样式，仿照 PDF 中"当事人信息"等居中带框标题
% 原档坐标：y≈539.1，字号约 15bp，方正大标宋，居中，有外框
% 与 notesection 同宽（467.72bp）、同X偏移对齐，紧贴上方 section 下方
% 用法：\subtitlesection{当事人信息}
\newcommand{\subtitlesection}[1]{%
    \begin{tcolorbox}[
        width=467.72bp,
        sharp corners,
        colback=white,
        colframe=inkblack,
        boxrule=0.5pt,
        boxsep=0pt,
        left=0pt, right=0pt,
        top=4bp, bottom=4bp,
        halign=center,
        fontupper=\fzdbs\fontsize{15bp}{15bp}\selectfont\color{inkblack},
        enlarge left by=-19.28bp,
        before={\nointerlineskip\vspace{-0.5pt}\noindent},
        after=\vspace{0pt}
    ]
    #1
    \end{tcolorbox}%
}

% 定义左右分栏表单区块，仿照 PDF 中"自诉人"等表单布局
% 左栏：标签垂直居中，右栏：表单字段，竖向分隔线
% #1: 左栏标签文字，#2: 右栏表单内容
% 用法：\pairsection{自诉人}{表单内容}
\newcommand{\pairsection}[2]{%
    \begin{tcolorbox}[
        width=467.72bp,
        sharp corners,
        colback=white,
        colframe=inkblack,
        boxrule=0.5pt,
        enlarge left by=-19.28bp,
        before={\nointerlineskip\vspace{-0.5pt}\noindent},
        after=\par,
        boxsep=0pt,
        left=0pt, right=0pt,
        top=0pt, bottom=0pt
    ]
    % 内部用 minipage 实现左右分栏
    \begin{minipage}[c]{112.9bp}%
        \centering\songti\fontsize{10.5bp}{10.5bp}\selectfont #1%
    \end{minipage}%
    \hspace{0pt}%
    \vrule width 0.5pt%
    \hspace{4bp}%
    \begin{minipage}[c]{346.32bp}%
        \vspace{4bp}%
        \songti\mdseries\fontsize{10.5bp}{17bp}\selectfont%
        \baselineskip=17bp%
        \setlength{\parindent}{0pt}%
        \color{inkblack}%
        #2%
        \vspace{4bp}%
    \end{minipage}%
    \end{tcolorbox}%
}

\begin{document}

% 严格还原 PDF 第一页的标题结构
\lawtitle{刑事（附带民事）自诉状}{（侮辱案）}

% 插入一个精确位置和宽度匹配的表框（根据原档：上Y=180.48，左X=70.86）
\begin{notesection}
\hfuzz=50pt % 忽略因中文字体度量导致的虚假宽度警告，绝对排版坐标已手工对齐完美边界
\color{inkblack}\heiti\fontsize{10.5bp}{16bp}\selectfont 说明：\par
\vspace*{-0.2bp}
\songti\mdseries\fontsize{10.5bp}{16bp}\selectfont % Fandol宋体，中等字重
\setlength{\parindent}{20.7bp} % 首行缩进两字符
\sloppy % 防止由于两端对齐导致的标点符号溢出越界

\setlength{\parskip}{0pt}
\vspace*{-5.1bp}为了方便您更好地参加诉讼，保护您的合法权利，请填写本表。
\lawpar
1.起诉时需提供自诉人、被告人的身份证或户口本等身份材料，如无法提供被告人的身份材料，\\[-4.86bp]
需提供被告人的联系电话、住址等信息。
\lawpar
2.委托律师为诉讼代理人的，需提供授权委托书、律师事务所公函及律师证复印件。\\[-4.86bp]
\hspace*{31.2bp}委托其他自然人为诉讼代理人的，需提供授权委托书、受托人身份材料、与自诉人关系的证\\[-4.86bp]
明等。
\lawpar
3.如果被害人死亡、丧失行为能力或者因受强制、威吓等无法告诉，或者是限制行为能力人以\\[-4.86bp]
及因年老、患病、盲、聋、哑等不能亲自告诉，其法定代理人、近亲属告诉或者代为告诉的，应当\\[-4.86bp]
提供与被害人关系的证明和被害人不能亲自告诉的原因的证明。
\lawpar
4.证据材料需说明证据名称和来源，有证人的，需提供证人的姓名、住址、联系方式等；有多\\[-4.86bp]
名证人的，请提供证人名单。
\lawpar
5.自诉人由于被告人的犯罪行为而遭受物质损失，同时提起附带民事诉讼的，需填写本表附带\\[-4.86bp]
民事部分有关内容。
\lawpar
6.本表有些内容可能与您的案件无关，您认为与案件无关的项目可以填“无”或不填；对于本\\[-4.86bp]
表中勾选项可以在对应项打“√”；您认为另有重要内容需要列明的，可以另附页填写。
\lawpar
7.本表word 电子版填写时，相关栏目可复制粘贴或扩容，但不得改变要素内容、格式设置。例\\[-4.86bp]
如，多自诉人、多被告人或多诉讼代理人等情况，可根据实际情况复制粘贴；需填写文字较多时，\\[-4.86bp]
可根据实际对栏目进行扩容等。
\lawpar
$\bigstar$特别提示$\bigstar$\\[-4.86bp]
\hspace*{20.7bp}如果诉讼参与人违反有关规定，虚假诉讼、恶意诉讼、滥用诉权，人民法院将视违法情形依法\\[-4.86bp]
追究责任。
\end{notesection}

% 当事人信息副标题
\subtitlesection{当事人信息}

% 自诉人表单
\pairsection{自诉人}{
姓名：\\
性别：男$\square$\hspace{1em}女$\square$\\
出生日期：\hspace{2em}年\hspace{2em}月\hspace{2em}日\\
民族：\\
出生地：\\
文化程度：\\
\makebox[0.5\linewidth][l]{职业：}\makebox[0.5\linewidth][l]{工作单位：}\\
户籍地：\\
住址：\\
联系电话：\\
证件类型：\\
证件号码：
}

\newpage

% ===== 第 2 页 =====

% 诉讼代理人
\pairsection{诉讼代理人}{
有$\square$\\
\hspace*{21bp}姓名：\\
\hspace*{21bp}\makebox[0.33\linewidth][l]{单位：}\makebox[0.33\linewidth][l]{职务：}\makebox[0.33\linewidth][l]{联系电话：}\\
\hspace*{21bp}（诉讼代理人为非律师的自然人，请增加填写以下信息）\\
\hspace*{21bp}住址：\\
\hspace*{21bp}证件类型：\\
\hspace*{21bp}证件号码：\\
\hspace*{21bp}与自诉人的关系：\\
无$\square$
}

% 法定代理人或代为告诉人
\pairsection{法定代理人或\\代为告诉人}{
有$\square$\\
\hspace*{21bp}姓名：\\
\hspace*{21bp}性别：男$\square$\hspace{1em}女$\square$\\
\hspace*{21bp}出生日期：\\
\hspace*{21bp}民族：\\
\hspace*{21bp}文化程度：\\
\hspace*{21bp}\makebox[0.5\linewidth][l]{职业：}\makebox[0.5\linewidth][l]{工作单位：}\\
\hspace*{21bp}住址：\\
\hspace*{21bp}联系电话：\\
\hspace*{21bp}证件类型：\\
\hspace*{21bp}证件号码：\\
\hspace*{21bp}与自诉人的关系：\\
无$\square$
}

% 被告人
\pairsection{被告人}{
姓名：\\
性别：男$\square$\hspace{1em}女$\square$\\
出生日期：\hspace{2em}年\hspace{2em}月\hspace{2em}日\\
民族：\\
出生地：\\
文化程度：\\
\makebox[0.5\linewidth][l]{职业：}\makebox[0.5\linewidth][l]{工作单位：}\\
户籍地：\\
住址：\\
联系电话：\\
证件类型：\\
证件号码：\\
发布侮辱信息的网络平台的名称及账号：
}

% 是否提起附带民事诉讼（单行条目，不使用 pairsection 以避免多余 padding）
\begin{tcolorbox}[
    width=467.72bp,
    sharp corners,
    colback=white,
    colframe=inkblack,
    boxrule=0.5pt,
    enlarge left by=-19.28bp,
    before={\nointerlineskip\vspace{-0.5pt}\noindent},
    after=\par,
    boxsep=0pt,
    left=0pt, right=0pt,
    top= 0pt, bottom=0pt
]
\begin{minipage}[c]{112.9bp}%
    \vspace{4pt}%
    \centering\songti\fontsize{10.5bp}{10.5bp}\selectfont 是否提起附带民事诉讼%
    \vspace{4pt}%
\end{minipage}%
\hspace{0pt}%
\vrule width 0.5pt%
\hspace{4bp}%
\begin{minipage}[c]{346.32bp}%
    \vspace{0.5pt}%
    \songti\mdseries\fontsize{10.5bp}{17bp}\selectfont%
    \setlength{\parindent}{0pt}%
    \color{inkblack}%
    \makebox[0.5\linewidth][l]{是$\square$}\makebox[0.5\linewidth][l]{否$\square$}%
    \vspace{0.5pt}%
\end{minipage}%
\end{tcolorbox}

\newpage

% ===== 第 3 页 =====

% 诉讼请求副标题
\subtitlesection{诉讼请求}

% 诉讼请求内容框
\begin{tcolorbox}[
    width=467.72bp,
    sharp corners,
    colback=white,
    colframe=inkblack,
    boxrule=0.5pt,
    enlarge left by=-19.28bp,
    before={\nointerlineskip\vspace{-0.5pt}\noindent},
    after=\par,
    boxsep=0pt,
    left=4bp, right=0pt,
    top=4bp, bottom=4bp
]
\songti\mdseries\fontsize{10.5bp}{17bp}\selectfont%
\baselineskip=17bp%
\parskip=0pt%
\lineskiplimit=-\maxdimen%
\setlength{\parindent}{0pt}%
\color{inkblack}%
1. 请求对被告人×××\hspace{0.5em}以侮辱罪追究刑事责任。\\
2.（提起附带民事诉讼的）请求被告人×××\hspace{0.5em}赔偿因犯罪行为给自诉人造成的物质损失。\\
3.（其他请求）。\\
通过信息网络实施侮辱行为，自诉人提供证据确有困难的，是否需要公安机关提供协助\\
是$\square$（具体事项和线索）\rule{120bp}{0.5pt}\\
否$\square$
\end{tcolorbox}

% 事实与理由副标题
\subtitlesection{事实与理由}

% 事实与理由内容框
\begin{tcolorbox}[
    width=467.72bp,
    sharp corners,
    colback=white,
    colframe=inkblack,
    boxrule=0.5pt,
    enlarge left by=-19.28bp,
    before={\nointerlineskip\vspace{-0.5pt}\noindent},
    after=\par,
    boxsep=0pt,
    left=4bp, right=0pt,
    top=4bp, bottom=4bp
]
\songti\mdseries\fontsize{10.5bp}{17bp}\selectfont%
\baselineskip=17bp%
\parskip=0pt%
\lineskiplimit=-\maxdimen%
\setlength{\parindent}{0pt}%
\color{inkblack}%
1. 事实（被告人实施侮辱行为的时间、地点、手段、情节、危害后果等，被告人通过网络实施的，自诉人如果与网络平台存在相关诉讼，请一并写明诉讼情况）：\\[51bp]
2. 理由（被告人涉嫌犯罪、承担附带民事赔偿责任的法律依据）：\\[51bp]
\mbox{}
\end{tcolorbox}

% 证据清单副标题
\begin{tcolorbox}[
    width=467.72bp,
    sharp corners,
    colback=white,
    colframe=inkblack,
    boxrule=0.5pt,
    boxsep=0pt,
    left=0pt, right=0pt,
    top=4bp, bottom=4bp,
    halign=center,
    enlarge left by=-19.28bp,
    before={\nointerlineskip\vspace{-0.5pt}\noindent},
    after=\vspace{0pt}
]
{\fzdbs\fontsize{15bp}{15bp}\selectfont\color{inkblack}证据清单}\par\nointerlineskip
\vspace{2bp}
{\heiti\fontsize{10.5bp}{10.5bp}\selectfont\color{inkblack}（证据材料另附）}
\end{tcolorbox}

% 证据清单内容框
\begin{tcolorbox}[
    width=467.72bp,
    sharp corners,
    colback=white,
    colframe=inkblack,
    boxrule=0.5pt,
    enlarge left by=-19.28bp,
    before={\nointerlineskip\vspace{-0.5pt}\noindent},
    after=\par,
    boxsep=0pt,
    left=4bp, right=0pt,
    top=4bp, bottom=4bp
]
\songti\mdseries\fontsize{10.5bp}{17bp}\selectfont%
\baselineskip=17bp%
\parskip=0pt%
\lineskiplimit=-\maxdimen%
\setlength{\parindent}{0pt}%
\color{inkblack}%
1.证明被告人实施侮辱行为、构成犯罪等证据材料\\
2.（提起附带民事诉讼的）证明因被告人实施侮辱行为给自诉人造成物质损失的证据材料\\
3.其他证据材料
\end{tcolorbox}

% 是否同意调解副标题
\subtitlesection{是否同意调解}

% 自诉部分
\begin{tcolorbox}[
    width=467.72bp,
    sharp corners,
    colback=white,
    colframe=inkblack,
    boxrule=0.5pt,
    enlarge left by=-19.28bp,
    before={\nointerlineskip\vspace{-0.5pt}\noindent},
    after=\par,
    boxsep=0pt,
    left=0pt, right=0pt,
    top=0pt, bottom=0pt
]
\begin{minipage}[c]{112.9bp}%
    \vspace{4pt}%
    \centering\songti\fontsize{10.5bp}{10.5bp}\selectfont 自诉部分%
    \vspace{4pt}%
\end{minipage}%
\hspace{0pt}%
\vrule width 0.5pt%
\hspace{4bp}%
\begin{minipage}[c]{346.32bp}%
    \vspace{0.5pt}%
    \songti\mdseries\fontsize{10.5bp}{17bp}\selectfont%
    \setlength{\parindent}{0pt}%
    \color{inkblack}%
    \makebox[0.33\linewidth][c]{同意$\square$}\makebox[0.33\linewidth][c]{不同意$\square$}\makebox[0.33\linewidth][c]{暂不确定$\square$}%
    \vspace{0.5pt}%
\end{minipage}%
\end{tcolorbox}

% 附带民事部分
\begin{tcolorbox}[
    width=467.72bp,
    sharp corners,
    colback=white,
    colframe=inkblack,
    boxrule=0.5pt,
    enlarge left by=-19.28bp,
    before={\nointerlineskip\vspace{-0.5pt}\noindent},
    after=\par,
    boxsep=0pt,
    left=0pt, right=0pt,
    top=0pt, bottom=0pt
]
\begin{minipage}[c]{112.9bp}%
    \vspace{4pt}%
    \centering\songti\fontsize{10.5bp}{10.5bp}\selectfont 附带民事部分%
    \vspace{4pt}%
\end{minipage}%
\hspace{0pt}%
\vrule width 0.5pt%
\hspace{4bp}%
\begin{minipage}[c]{346.32bp}%
    \vspace{0.5pt}%
    \songti\mdseries\fontsize{10.5bp}{17bp}\selectfont%
    \setlength{\parindent}{0pt}%
    \color{inkblack}%
    \makebox[0.33\linewidth][c]{同意$\square$}\makebox[0.33\linewidth][c]{不同意$\square$}\makebox[0.33\linewidth][c]{暂不确定$\square$}%
    \vspace{0.5pt}%
\end{minipage}%
\end{tcolorbox}

% 具状人签字与日期
\vspace{10bp}
\noindent\hspace*{248bp}{\fzdbs\fontsize{15bp}{22bp}\selectfont 具状人（签字）：}\par
\noindent\hspace*{248bp}{\fzdbs\fontsize{15bp}{22bp}\selectfont 日期：}

\end{document}
